\documentclass[11pt]{article}
\usepackage[margin=1in]{geometry}
\usepackage{booktabs}
\usepackage{hyperref}
\usepackage{amsmath}
\usepackage{graphicx}

\title{Negative Feedback and Reach in Short-Video Recommendation:\\
A Deterministic OASIS Study}
\author{ZiyangLiu\\\texttt{me.zyliu@gmail.com}}
\date{April 2026}

\begin{document}
\maketitle

\begin{abstract}
We study whether explicit negative feedback can reduce the reach of low-quality
content in a short-video recommender. Using OASIS, we first run a deterministic
two-condition TikTok-style simulation (baseline versus stronger
\texttt{not\_interested} treatment) to verify mechanism behavior: platform-level
totals remain stable while target-item negative feedback increases and target
traffic-pool level decreases. We then run a larger four-arm robustness study
(18 parameter configurations, 30 paired runs each, with scaled
creators/viewers/videos). At this larger scale, treatment robustly increases
target negative-feedback metrics (18/18 significant), while target pool-level
demotion remains directionally negative but statistically weaker (3/18
significant). The findings support explicit negative feedback as a meaningful
governance signal, with effect strength depending on regime and scale.
\end{abstract}

\section{Introduction}
Negative-feedback controls such as ``not interested'' are widely deployed in
short-video products, but their direct effect is difficult to isolate in live
systems. Online traffic is noisy, policies change over time, and experiments
must balance rigor against operational risk.

We use agent-based simulation to study this mechanism under controlled
conditions. OASIS~\cite{oasis2024} supports replayable social experiments with
configurable users, actions, and ranking behavior. The scope here is narrow:
whether increasing explicit negative feedback on one low-quality item changes
its distribution level while coarse platform totals remain stable.

\paragraph{Contributions}
\begin{enumerate}
\item A deterministic two-condition protocol for short-video feedback analysis.
\item A controlled result showing target-item demotion (\texttt{traffic\_pool\_level}: 1 $\rightarrow$ 0).
\item Reproducible artifacts (DBs, reports, and comparison outputs) for extension.
\end{enumerate}

\section{Related Work}
Recent short-video recommendation work has emphasized modeling passive-negative
feedback from user skip behavior and integrating such signals in sequential
ranking objectives (\cite{pan2023passivenegative}). That line of work is
primarily offline and model-centric, evaluated on logged behavioral datasets.

Our study is complementary rather than competitive with that objective. We focus
on policy-effect evaluation in a controlled simulation: given matched
conditions, what changes in distribution emerge when explicit negative feedback
is intensified for a target low-quality item? In other words, this paper asks a
mechanism and intervention question, while prior work mainly asks a prediction
and ranking-accuracy question.

\section{Method}
\subsection{Environment}
We use the TikTok-style mode in OASIS with traffic-pool recommendation logic.
The run contains 3 creators, 6 viewers, and 4 uploaded videos (dance, tech,
comedy, lifestyle). Agents perform watches, comments, shares, and negative
feedback (\texttt{not\_interested}). Livestream actions are enabled but are not
primary outcomes in this paper.

Our ranking setup follows standard industrial recommendation practice where
engagement and satisfaction signals are combined at ranking time
(\cite{covington2016youtube,zhao2019multitask}) and where feedback signals can
be interpreted as preference evidence (\cite{rendle2009bpr}).

\subsection{Experimental Design}
We construct two deterministic conditions:
\begin{itemize}
\item \textbf{Baseline:} weak negative feedback against one low-quality comedy video.
\item \textbf{Treatment:} stronger negative feedback against the same target video.
\end{itemize}
All non-treatment factors are identical: user identities, follow graph, upload
sequence, viewer plans, ranking parameters, and simulation step progression. The
only intervention is additional negative-feedback events on the comedy item in
treatment.

\subsection{Metrics}
Platform-level:
\begin{itemize}
\item \texttt{total\_views}
\item \texttt{avg\_watch\_ratio}
\end{itemize}
Target-video level (low-quality comedy):
\begin{itemize}
\item \texttt{negative\_count}
\item \texttt{traffic\_pool\_level}
\end{itemize}

\section{Experimental Setup and Reproducibility}
The pipeline consists of:
\begin{enumerate}
\item Generating baseline/treatment databases via
\texttt{examples/experiment/short\_video\_negative\_feedback/run\_experiment.py}.
\item Generating per-run summaries via
\texttt{visualization/short\_video\_simulation/code/generate\_report.py}.
\item Generating cross-run comparison via
\texttt{visualization/short\_video\_simulation/code/compare\_runs.py}.
\end{enumerate}

Key artifacts include:
\texttt{negative\_feedback\_baseline.db},
\texttt{negative\_feedback\_treatment.db},
\texttt{baseline\_report.md},
\texttt{treatment\_report.md},
\texttt{comparison\_time\_series.csv}, and
\texttt{comparison\_metrics.png}.

\section{Results}
\subsection{Main Aggregate Results (Single-Run)}
\begin{table}[h]
\centering
\begin{tabular}{lrr}
\toprule
Metric & Baseline & Treatment \\
\midrule
total\_views & 16 & 16 \\
avg\_watch\_ratio & 0.638 & 0.638 \\
target negative\_count & 1 & 3 \\
target traffic\_pool\_level & 1 & 0 \\
\bottomrule
\end{tabular}
\caption{Main outcomes under matched baseline and treatment conditions.}
\label{tab:main-results}
\end{table}

Platform-level outcomes remain stable while the target low-quality content is
demoted under stronger negative feedback.

\subsection{Time-Series Evidence}
Figure~\ref{fig:comparison-metrics} reports the run-level trajectories generated
from the comparison artifact (\texttt{comparison\_metrics.png}), including
cumulative views, average watch ratio, and average traffic-pool level.

\begin{figure}[h]
\centering
\includegraphics[width=\linewidth]{../../../../examples/experiment/short_video_negative_feedback/output/comparison/comparison_metrics.png}
\caption{Baseline vs. treatment trajectories for cumulative views, average
watch ratio, and average traffic-pool level.}
\label{fig:comparison-metrics}
\end{figure}

The key divergence appears at the low-quality comedy-item update step, where the
treatment condition accumulates stronger negative feedback and triggers item
demotion. Outside the treated item, global curves remain close, consistent with
a selective intervention effect.

\subsection{Multi-Run Statistical Results ($n=30$)}
To increase robustness, we ran 30 paired baseline/treatment experiments with
seeded perturbations
(\texttt{watch\_jitter}=0.08, \texttt{behavior\_flip\_prob}=0.05,
\texttt{treatment\_extra\_neg\_prob}=0.85). Table~\ref{tab:multirun-results}
reports paired differences and two-sided paired sign-flip p-values.

\begin{table}[h]
\centering
\begin{tabular}{lrrr}
\toprule
Metric & Paired Diff Mean & 95\% CI & p-value \\
\midrule
total\_negative\_feedback & +1.700 & [1.372, 2.028] & 0.00005 \\
comedy\_negative\_count & +1.767 & [1.613, 1.921] & 0.00005 \\
comedy\_traffic\_pool\_level & -0.667 & [-0.862, -0.471] & 0.00005 \\
avg\_watch\_ratio & -0.00068 & [-0.00445, 0.00308] & 0.72296 \\
total\_shares & -0.100 & [-0.355, 0.155] & 0.60457 \\
\bottomrule
\end{tabular}
\caption{Paired multi-run treatment effects (Treatment - Baseline), 30 runs.}
\label{tab:multirun-results}
\end{table}

These multi-run results strengthen the mechanism claim: treatment consistently
increases negative feedback and lowers the target comedy video's traffic-pool
level, while global engagement-style metrics do not show statistically
meaningful shifts in this setup.

\subsection{Large-Scale Parameter Sweep (18 Configurations $\times$ 30 Runs, Four Arms)}
We strengthen the evaluation with a four-arm setup (\texttt{baseline},
\texttt{treatment}, \texttt{none}, \texttt{random}) and larger synthetic scale
(3x creators, 5x viewers, 4x videos). For the treatment-versus-baseline
comparison, we keep the same 18-configuration parameter grid:
\texttt{watch\_jitter} $\in \{0.04,0.08,0.12\}$,
\texttt{behavior\_flip\_prob} $\in \{0.02,0.05\}$, and
\texttt{treatment\_extra\_neg\_prob} $\in \{0.70,0.85,1.00\}$.

Under this stronger setup, \texttt{comedy\_negative\_diff\_mean} remains
positive and significant in 18/18 configurations, while
\texttt{comedy\_pool\_diff\_mean} is directionally negative on average but only
significant in 3/18 configurations. Platform-level \texttt{watch\_ratio} and
\texttt{shares} remain non-significant in 18/18 configurations.

\begin{table}[h]
\centering
\begin{tabular}{lrrrrrr}
\toprule
\texttt{treatment\_extra\_neg\_prob} & Pool Mean & Pool SD & Pool Sig & Neg Mean & Neg SD & Neg Sig \\
\midrule
0.70 & -0.061 & 0.039 & 0/6 & +1.339 & 0.071 & 6/6 \\
0.85 & -0.083 & 0.072 & 1/6 & +1.567 & 0.060 & 6/6 \\
1.00 & -0.117 & 0.084 & 2/6 & +1.922 & 0.046 & 6/6 \\
\bottomrule
\end{tabular}
\caption{Large-scale treatment robustness grouped by treatment intensity.}
\label{tab:batch-tn-summary}
\end{table}

The aggregate trend remains monotonic: stronger treatment intensity increases
target negative-feedback accumulation and yields more negative (though modest)
target pool-level shifts.

\subsection{Control-Arms Comparison (vs Baseline)}
To test whether any perturbation can mimic treatment effects, we compare
\texttt{none}, \texttt{random}, and \texttt{treatment} against baseline across
the same 18 configurations.

\begin{table}[h]
\centering
\begin{tabular}{lrrrr}
\toprule
Condition & Comedy Neg Diff (mean $\pm$ sd) & Comedy Pool Diff (mean $\pm$ sd) & Watch Diff (mean $\pm$ sd) & Share Diff (mean $\pm$ sd) \\
\midrule
none & -5.339 $\pm$ 0.169 & +0.022 $\pm$ 0.040 & +0.000000 $\pm$ 0.000000 & +0.000 $\pm$ 0.000 \\
random & +0.700 $\pm$ 0.307 & +0.011 $\pm$ 0.044 & -0.000063 $\pm$ 0.000766 & +0.085 $\pm$ 0.724 \\
treatment & +1.609 $\pm$ 0.253 & -0.087 $\pm$ 0.068 & +0.000090 $\pm$ 0.000429 & +0.026 $\pm$ 0.351 \\
\bottomrule
\end{tabular}
\caption{Control-arm contrasts versus baseline, aggregated over 18 configurations.}
\label{tab:control-arms}
\end{table}

Random negatives weakly increase target negative-count but do not induce
targeted demotion. Treatment remains the only arm with a consistent negative
direction on target pool-level difference.

\begin{figure}[h]
\centering
\includegraphics[width=\linewidth]{../../../../examples/experiment/short_video_negative_feedback/batch_output_85plus/batch_85plus_fourarm_summary.png}
\caption{Large-scale four-arm robustness summary (18 configurations, 30 paired runs each).}
\label{fig:batch-robustness}
\end{figure}

\subsection{Mechanism Calibration: Balanced-Sensitive Profile ($n=30$)}
To improve mechanism observability in the scaled setting, we introduced a
\texttt{balanced\_sensitive} recommender profile that increases
negative-feedback penalty while avoiding floor effects from overly aggressive
demotion. We then ran a paired 30-run probe (20x creators/viewers/videos,
\texttt{watch\_jitter}=0.30, \texttt{behavior\_flip\_prob}=0.10,
\texttt{treatment\_extra\_neg\_prob}=1.00).

\begin{table}[h]
\centering
\begin{tabular}{lrrrr}
\toprule
Metric & Paired Diff Mean & p-value & Baseline Mean & Treatment Mean \\
\midrule
negative\_feedback\_rate & +0.005646 & 0.000050 & 0.151547 & 0.157193 \\
comedy\_traffic\_pool\_level & -0.233333 & 0.015249 & 0.633333 & 0.400000 \\
retention\_3s\_rate & +0.000000 & 1.000000 & 0.982271 & 0.982271 \\
avg\_watch\_ratio & +0.000000 & 1.000000 & 0.622830 & 0.622830 \\
creator\_coverage & +0.000000 & 1.000000 & 1.000000 & 1.000000 \\
\bottomrule
\end{tabular}
\caption{Balanced-sensitive calibration probe (Treatment - Baseline), 30 paired runs.}
\label{tab:balanced-calibration}
\end{table}

This calibration closes an important gap from earlier large-scale runs:
targeted negative feedback no longer changes only count-level signals; it also
induces statistically significant target demotion at traffic-pool level.

\section{Discussion}
The large-scale result is informative and more nuanced. Targeted explicit
negative feedback robustly increases item-level negative signals, while
pool-level demotion remains sensitive to parameter regime under broader noise
and scale. With calibrated ranking sensitivity (\texttt{balanced\_sensitive}),
targeted demotion becomes statistically significant while platform-level watch
and coverage signals stay stable. This suggests that negative feedback is a
strong governance signal, but its distributional effect depends on calibration
of ranking sensitivity.

\section{Limitations}
\begin{enumerate}
\item The simulation world is still small (fixed creators, viewers, and categories) relative to real platform scale.
\item Short horizon and limited content taxonomy.
\item No strategic adaptation from creators or viewers.
\item Findings show mechanism behavior, not production-scale effect sizes.
\end{enumerate}

\section{Ethics and Broader Impacts}
This study uses synthetic agents and no personal user data. In real deployment,
negative-feedback-driven demotion can still raise fairness and representation
concerns. Any deployment should pair these signals with calibration, monitoring,
and appeals-aware processes.

\section{Conclusion}
Under controlled conditions in OASIS, increasing explicit negative feedback on a
low-quality video reliably increases target negative-feedback metrics while
keeping platform-level totals stable. With the calibrated
\texttt{balanced\_sensitive} profile, we also observe statistically significant
target pool-level demotion in the scaled setting
(\texttt{comedy\_traffic\_pool\_level} diff $-0.233333$, $p=0.015249$). This
provides a stronger mechanism baseline for follow-up moderation and
ranking-policy studies.

This use of synthetic populations is also aligned with recent simulation-driven
behavioral experimentation in agent systems (\cite{park2023generative}).

\section{Data and Reproducibility}
All core artifacts are available in the repository output directory:
\begin{itemize}
\item \texttt{examples/experiment/short\_video\_negative\_feedback/output/negative\_feedback\_baseline.db}
\item \texttt{examples/experiment/short\_video\_negative\_feedback/output/negative\_feedback\_treatment.db}
\item \texttt{examples/experiment/short\_video\_negative\_feedback/output/comparison/comparison\_time\_series.csv}
\item \texttt{examples/experiment/short\_video\_negative\_feedback/output/comparison/comparison\_metrics.png}
\item \texttt{examples/experiment/short\_video\_negative\_feedback/batch\_output\_85plus/batch\_config\_summary.csv}
\item \texttt{examples/experiment/short\_video\_negative\_feedback/batch\_output\_85plus/batch\_report.md}
\item \texttt{examples/experiment/short\_video\_negative\_feedback/batch\_output\_85plus/batch\_85plus\_fourarm\_summary.png}
\item \texttt{examples/experiment/short\_video\_negative\_feedback/output\_metric\_probe\_balanced\_n30/multirun\_summary.csv}
\end{itemize}

Reproduction entry points:
\begin{itemize}
\item \texttt{examples/experiment/short\_video\_negative\_feedback/run\_experiment.py}
\item \texttt{examples/experiment/short\_video\_negative\_feedback/run\_batch\_experiments.py}
\item \texttt{visualization/short\_video\_simulation/code/generate\_report.py}
\item \texttt{visualization/short\_video\_simulation/code/compare\_runs.py}
\item \texttt{scripts/ci/run\_short\_video\_metric\_probe\_balanced\_n30.sh}
\end{itemize}

\section*{Acknowledgments}
The author thanks the OASIS project and community for the open-source framework
and documentation that made this study possible.

\bibliographystyle{plain}
\bibliography{references}

\end{document}
